\chapter{Marco teórico y estado del arte}

\label{Capitulo_2}

%----------------------------------------------------------------------------------------
\section{Introducción}
\label{sec:estado_arte_intro}

En este capítulo se presentan los fundamentos teóricos de los métodos de
resolución utilizados en esta tesina y se revisan los trabajos relacionados
con el problema de ruteo de picking.

%----------------------------------------------------------------------------------------
\section{Modelo de Programación Lineal Entera Mixta}
\label{sec:marco_mip}

% TODO: completar

%----------------------------------------------------------------------------------------
\section{Hill Climbing}
\label{sec:marco_hc}

El algoritmo de Hill Climbing, también referido frecuentemente como búsqueda local, descenso o mejora iterativa, es considerado el
método metaheurístico más antiguo y sencillo. Pertenece a la familia de las metaheurísticas basadas en una sola solución, cuyo
funcionamiento principal consiste en mejorar una solución individual realizando recorridos o trayectorias a través de las vecindades
del espacio de búsqueda de un problema \cite{talbi2009metaheuristics}.

El concepto central de toda metaheurística basada en una sola solución es el de \emph{vecindad}. Formalmente, una función de
vecindad $N : S \rightarrow 2^S$ asigna a cada solución $s$ del espacio de búsqueda $S$ un conjunto de soluciones $N(s) \subset S$,
llamadas \emph{vecinas} de $s$. Cada vecino se obtiene aplicando sobre $s$ un \emph{operador de movimiento}, que realiza una pequeña
perturbación de la solución. A partir de esta definición, una solución $s$ es un \emph{óptimo local} respecto de $N$ si
$f(s) \leq f(s')$ para todo $s' \in N(s)$, en un problema de minimización. Una solución puede ser un óptimo local para una vecindad
y no serlo para otra, por lo que la elección de la vecindad determina en gran medida el comportamiento de la búsqueda
\cite{talbi2009metaheuristics}.

En problemas cuya solución se representa como una permutación, como el problema del viajante (TSP, del inglés \emph{Traveling
Salesman Problem}), los operadores de movimiento más utilizados son el \emph{intercambio} (\emph{swap}), que intercambia la posición
de dos elementos y genera una vecindad de tamaño $n(n-1)/2$; la \emph{inserción}, que extrae un elemento de su posición y lo coloca
en otra; y el operador \emph{k-opt}, que elimina $k$ aristas del recorrido y las reemplaza por otras $k$. El caso más difundido de
este último es \emph{2-opt}, propuesto por Croes \cite{croes1958method}, que reemplaza las aristas $(\pi_i, \pi_{i+1})$ y $(\pi_j, \pi_{j+1})$ por $(\pi_i, \pi_j)$ y
$(\pi_{i+1}, \pi_{j+1})$. Esto equivale a invertir el tramo del recorrido comprendido entre ambas aristas:
\begin{equation}
	\pi'(k) =
	\begin{cases}
		\pi(k)         & \text{si } k \leq i \text{ o } k > j, \\
		\pi(i + j + 1 - k) & \text{si } i < k \leq j.
	\end{cases}
	\label{eq:2opt}
\end{equation}
La vecindad de 2-opt tiene $n(n-1)/2 - n$ elementos, ya que se consideran todos los pares de aristas excepto los adyacentes
\cite{talbi2009metaheuristics}.

El proceso de optimización comienza siempre evaluando una solución inicial dada. A partir de ese punto, y durante cada iteración, el
algoritmo reemplaza la solución actual por un vecino que logre mejorar el valor de la función objetivo. Para decidir qué camino tomar
en el espacio de búsqueda, el método requiere definir una estrategia de selección del próximo vecino, existiendo opciones estándar como
elegir la mejor mejora disponible, tomar la primera mejora que el algoritmo encuentre, o seleccionar una mejora de forma aleatoria
\cite{talbi2009metaheuristics}. La estrategia de primera mejora aplica pasos de mejora antes que la de mejor mejora, aunque cada paso
suele mejorar menos y se necesitan más pasos para alcanzar un óptimo local. Además, el óptimo local que encuentra depende del orden en
que se recorre la vecindad: recorrerla en orden aleatorio permite que distintas ejecuciones, incluso partiendo de la misma solución
inicial, alcancen óptimos locales muy diferentes \cite{hoos2015sls}.

Una característica particular del Hill Climbing es que no requiere configurar condiciones de finalización complejas, ya que sus
criterios de parada son implícitos. La búsqueda se detiene de forma natural cuando todos los vecinos candidatos evaluados resultan ser
peores que la solución actual. Cuando esto sucede, significa que el algoritmo ha encontrado un óptimo local y concluye su ejecución
\cite{talbi2009metaheuristics}. Sin embargo, para confirmar que una solución es un óptimo local es necesario inspeccionar su vecindad
completa \cite{hoos2015sls}. Cuando en cada iteración se genera solo una parte de la vecindad, como ocurre al muestrear vecinos al
azar, la búsqueda se detiene mediante criterios explícitos, como un número máximo de iteraciones o un límite de tiempo
\cite{talbi2009metaheuristics}.

El principal inconveniente de esta metaheurística es precisamente su tendencia a quedar atrapada en dichos óptimos locales, limitando
su capacidad para encontrar el óptimo global del problema. Para contrarrestar esta debilidad, se han desarrollado cuatro grandes
familias de enfoques: reiniciar el algoritmo iterando desde diferentes soluciones iniciales, cambiar la estructura de la vecindad de
búsqueda, alterar la función objetivo o los datos de entrada para modificar el panorama del problema, o finalmente, flexibilizar el
criterio de aceptación para permitir movimientos hacia vecinos que no mejoren la solución actual, dando origen a técnicas más avanzadas
como el recocido simulado y la búsqueda tabú \cite{talbi2009metaheuristics, hoos2015sls}.

%----------------------------------------------------------------------------------------
\section{Ant Colony Optimization}
\label{sec:marco_aco}

La Optimización por Colonia de Hormigas, conocida por sus siglas en inglés como ACO (Ant Colony Optimization), es una
metaheurística poblacional introducida por M. Dorigo que se inspira en el comportamiento de búsqueda de alimento de las hormigas
reales y su capacidad para encontrar los caminos más cortos entre su nido y una fuente de comida. Este fenómeno biológico se basa
en la \emph{estigmergia}, una forma de comunicación indirecta en la que las hormigas modifican su entorno al depositar una
sustancia química olorosa llamada \emph{feromona} a lo largo de su trayectoria. En el ámbito computacional, este comportamiento
coordinado se modela mediante el uso de ``hormigas artificiales'', que actúan como agentes estocásticos que cooperan para
construir de manera incremental soluciones a problemas de optimización combinatoria complejos moviéndose a través de un grafo
de construcción \cite{dorigo2004aco}.

El proceso de búsqueda en ACO se basa en la construcción probabilística de soluciones \cite{hoos2015sls}. A medida que una hormiga
avanza por el grafo, toma decisiones sobre el siguiente paso basándose en la interacción de dos factores fundamentales: los rastros de
feromona artificial y la información heurística dependiente del problema. Así, las hormigas realizan una búsqueda constructiva
aleatorizada pero fuertemente sesgada hacia opciones que parecen prometedoras tanto a nivel heurístico como según la experiencia previa
de la colonia \cite{dorigo2004aco}. Por un lado, la feromona representa la memoria a largo plazo y compartida del proceso de búsqueda,
indicando qué componentes han formado parte de buenas soluciones en el pasado; por otro, la información heurística aporta conocimiento
previo sobre el problema específico, como por ejemplo priorizar las distancias más cortas entre dos puntos \cite{talbi2009metaheuristics}.

En el \emph{Ant System} (AS), el primer algoritmo ACO, esta interacción se formaliza mediante la \emph{regla proporcional aleatoria}:
la probabilidad de que la hormiga $k$, ubicada en el nodo $i$, elija moverse al nodo $j$ es
\begin{equation}
	p_{ij}^{k} = \frac{[\tau_{ij}]^{\alpha} \, [\eta_{ij}]^{\beta}}{\sum_{l \in \mathcal{N}_i^k} [\tau_{il}]^{\alpha} \, [\eta_{il}]^{\beta}},
	\qquad \text{si } j \in \mathcal{N}_i^k,
	\label{eq:aco_transicion}
\end{equation}
donde $\tau_{ij}$ es la feromona asociada al arco $(i, j)$; $\eta_{ij} = 1/d_{ij}$ es la información heurística, definida como la
inversa de la distancia entre ambos nodos; $\mathcal{N}_i^k$ es el conjunto de nodos que la hormiga $k$ todavía no visitó, y fuera
del cual la probabilidad de elección es nula; y $\alpha$ y $\beta$ son parámetros que regulan la influencia relativa de la feromona
y de la información heurística. Si $\alpha = 0$, se favorecen los nodos más cercanos y el algoritmo se comporta como una heurística
golosa estocástica. Si $\beta = 0$, solo actúa la feromona, lo que en general produce malos resultados y, para $\alpha > 1$, lleva
rápidamente a una situación de \emph{estancamiento}, en la que todas las hormigas construyen el mismo recorrido \cite{dorigo2004aco}.

Una vez que las hormigas han construido sus soluciones completas, el algoritmo procede a actualizar los niveles de feromona mediante dos
mecanismos opuestos y complementarios: la evaporación y el depósito o refuerzo, que incrementa la cantidad de feromona en aquellos
componentes que pertenecen a las soluciones generadas en la iteración actual \cite{hoos2015sls}. La evaporación reduce todos los niveles
de feromona por un factor constante a lo largo del tiempo. Este proceso de ``olvido'' resulta crucial para evitar una convergencia
prematura hacia soluciones subóptimas y fomentar la exploración de nuevas áreas del espacio de búsqueda. En cuanto al depósito, en la
mayoría de las variantes de alto rendimiento la cantidad de feromona depositada es proporcional a la calidad de la solución, y a menudo
solo se permite que las hormigas con las mejores rutas refuercen sus caminos, atrayendo así a las futuras hormigas hacia las regiones
más óptimas \cite{dorigo2004aco}.

En AS, ambos mecanismos se aplican una vez que las $m$ hormigas completaron sus recorridos. Primero, la evaporación reduce la
feromona de todos los arcos:
\begin{equation}
	\tau_{ij} \leftarrow (1 - \rho) \, \tau_{ij},
	\label{eq:aco_evaporacion}
\end{equation}
donde $0 < \rho \leq 1$ es la \emph{tasa de evaporación}. Si un arco deja de ser elegido, su feromona decrece exponencialmente con el
número de iteraciones. Luego, cada hormiga deposita feromona sobre los arcos que recorrió:
\begin{equation}
	\tau_{ij} \leftarrow \tau_{ij} + \sum_{k=1}^{m} \Delta\tau_{ij}^{k},
	\qquad
	\Delta\tau_{ij}^{k} =
	\begin{cases}
		1 / C^{k} & \text{si el arco } (i, j) \text{ pertenece a } T^{k}, \\
		0         & \text{en caso contrario,}
	\end{cases}
	\label{eq:aco_deposito}
\end{equation}
donde $C^{k}$ es la longitud del recorrido $T^{k}$ construido por la hormiga $k$. De esta forma, cuanto mejor es el recorrido de una
hormiga, más feromona reciben sus arcos \cite{dorigo2004aco}.

Sobre este esquema se desarrollaron variantes que explotan con más fuerza las mejores soluciones encontradas. En el \emph{Elitist Ant
System} (EAS), además del depósito de todas las hormigas, los arcos del mejor recorrido encontrado desde el inicio del algoritmo,
$T^{bs}$ (del inglés \emph{best-so-far}), reciben un refuerzo adicional $e / C^{bs}$, donde $e$ es un parámetro que define el peso de
ese recorrido y $C^{bs}$ su longitud. En el \emph{MAX-MIN Ant System} (MMAS), propuesto por Stützle y Hoos \cite{stutzle2000mmas}, en cambio, solo la mejor hormiga de la iteración o la
mejor desde el inicio deposita feromona, con $\Delta\tau_{ij}^{best} = 1 / C^{best}$. Como esta estrategia puede llevar al
estancamiento, MMAS además limita la feromona al intervalo $[\tau_{\min}, \tau_{\max}]$ y la inicializa en $\tau_{\max}$ para
favorecer la exploración al comienzo de la búsqueda \cite{dorigo2004aco}.

Para maximizar la calidad de los resultados, los algoritmos ACO suelen incorporar acciones centralizadas que no pueden realizar las hormigas a
nivel individual, conocidas como ``acciones de demonio''. Una de las estrategias más extendidas y recomendables es aplicar algoritmos de búsqueda
local o procedimientos de mejora iterativa a las soluciones completas antes de proceder a la actualización de las feromonas \cite{dorigo2004aco}.

Si bien esta metaheurística se diseñó inicialmente para problemas de un solo objetivo, su estructura se ha adaptado exitosamente para abordar
problemas de optimización multiobjetivo (MOP). En estas adaptaciones, el marco de trabajo adquiere una mayor complejidad e incluye variaciones
como el uso de múltiples colonias independientes o la implementación de varias estructuras de matrices de feromonas, donde cada rastro se asocia
a una función objetivo específica. Durante la construcción de la ruta, una hormiga puede seleccionar de forma aleatoria o ponderada a qué matriz
de feromona prestar atención, lo que permite al algoritmo converger hacia un conjunto de soluciones óptimas de Pareto que equilibran múltiples
metas contradictorias simultáneamente \cite{alaya2007antcolony}.

%----------------------------------------------------------------------------------------
\section{Manejo de restricciones en metaheurísticas}
\label{sec:marco_restricciones}

Un aspecto común al diseño de cualquier metaheurística, incluidas \ac{hc} y \ac{aco}, es el tratamiento de las restricciones del
problema. La estrategia más simple es la de \emph{rechazo}, que descarta automáticamente toda solución infactible y conserva solo las
factibles durante la búsqueda. Esta estrategia es razonable cuando la porción infactible del espacio de búsqueda es muy pequeña, pero
desaprovecha la información que aportan las soluciones infactibles, aun cuando las soluciones óptimas suelen encontrarse en la frontera
entre la región factible y la infactible. Además, cuando la región factible es discontinua, el único camino entre dos soluciones
factibles puede atravesar soluciones infactibles \cite{talbi2009metaheuristics}.

Por este motivo, el enfoque más difundido es la estrategia de \emph{penalización}, en la que las soluciones infactibles sí se
consideran durante la búsqueda, pero la función objetivo se extiende con un término que las penaliza. En su forma lineal, la
función penalizada es
\begin{equation}
	f'(s) = f(s) + \lambda \, c(s),
	\label{eq:penalizacion}
\end{equation}
donde $c(s)$ representa el costo de la violación de las restricciones y $\lambda$ es un peso. Esto permite que la búsqueda pase por
una solución infactible para llegar a una factible mejor que la anterior. La penalización resulta más informativa cuando $c(s)$ mide
\emph{cuánto} se viola cada restricción, y no solo \emph{cuántas} restricciones se violan, ya que así refleja qué tan cerca está la
solución de la región factible. La elección del peso es delicada: si es demasiado bajo, la solución final puede resultar
infactible; si es demasiado alto, la búsqueda puede converger hacia soluciones factibles pero no óptimas. El peso puede mantenerse
constante durante toda la búsqueda (penalización \emph{estática}) o modificarse a medida que avanza (penalización \emph{dinámica})
\cite{talbi2009metaheuristics}.

%----------------------------------------------------------------------------------------
\section{Trabajos relacionados}
\label{sec:trabajos_relacionados}

El modelo de referencia de esta tesina es el propuesto por Coruzzolo et al.
\cite{coruzzolo2023orderpicking}, que formula la optimización conjunta de
\emph{batching}, asignación de lotes, secuenciación y ruteo de picking como
un modelo de \ac{mip}, y lo valida sobre un caso real
reduciendo el tiempo de picking en un 57\% respecto de una estrategia de
pedido por pedido.

Otros trabajos documentan el impacto económico del picking y motivan porqué
vale la pena optimizarlo: Urzúa et al. \cite{urzua2019evaluating} evalúan,
mediante simulación de eventos discretos, el efecto de distintas estrategias
de picking (\emph{slotting}, \emph{wave picking}, \emph{picking} expedito)
sobre el tiempo de cumplimiento de pedidos, con reducciones de hasta el 54\%;
Gutiérrez López et al. \cite{gutierrez2019reduction} documentan un caso real
de reducción del tiempo de picking mediante la reubicación de productos en un
centro de distribución.

En cuanto a los métodos de resolución, Chen et al.
\cite{chen2013antcolony} proponen un algoritmo de ruteo basado en \ac{aco}
para dos operarios de picking que considera explícitamente la congestión
entre ellos, un fenómeno relevante para cualquier heurística de ruteo con
múltiples operarios trabajando en simultáneo. Lesch et al.
\cite{lesch2023optimizing} combinan un algoritmo genético para la asignación
de \emph{storage} con \ac{aco} para el ruteo de picking en almacenes con
entrepiso, mostrando que ambos problemas están interrelacionados y que
optimizarlos de forma conjunta mejora el resultado global. Alaya et al.
\cite{alaya2007antcolony} formalizan extensiones de \ac{aco} a problemas
multiobjetivo, relevantes porque el OPP combina en su función objetivo
tiempo de recolección con penalidades por adelanto y demora.
